We consider a population of $N$ language model (LM) agents that hold opinions about a shared question, communicate with each other, and revise their opinions through multiple rounds of interactions (Figure \ref{fig:1_intro}).

\textit{Personas and Questions.} The agents answer one of two kinds of questions. \emph{Objective questions} are competition mathematics problems from the MATH dataset \citep{hendrycks2021measuringmathematicalproblemsolving}, recast as binary multiple choice questions. For example, \textit{``Find the distance between the foci of the hyperbola
$x^2 - 6x - 4y^2 - 8y = 27$. A: $4\sqrt{5}$\quad B: $4\sqrt{10}$''}. Here an agent's persona encodes \emph{expertise}: each agent is shown the worked ground-truth solution to a different MATH problem, so agents differ in the expertise they bring to the shared question. \emph{Subjective questions} are political statements adapted from the Political Questions Dataset for LLM Bias Evaluation \citep{promptfoo_political_bias_2025}, which each agent either agrees ($+1$) or disagrees ($-1$) with. For example,
\textit{``Mandatory vaccination violates bodily autonomy''}. Here personas are short free-text profiles adapted from TWIN-2K-500 \citep{toubia2025twin2k500datasetbuildingdigital}, describing an individual's demographics, political alignment, and personality traits, so agents differ in preference. Appendix~\ref{apdx:tasks} gives the full construction of both datasets.

\textit{Agents.} Each agent $i$ carries a fixed persona $p_i$, which distinguishes it from other agents and introduces heterogeneity into the population. The persona reflects variation in personal preferences (subjective) and competence (objective).\footnote{For the agents answering subjective questions, personas are based on profiles of real individuals from TWIN-2K-500
\citep{toubia2025twin2k500datasetbuildingdigital}. For objective questions, we encode
$p_i$ as domain expertise using the ground truth answers from \citet{hendrycks2021measuringmathematicalproblemsolving}.} At each time step $t$, agent $i$ processes an input $x_i(t)$, which consists of messages the agent $i$ received from its connections, and samples an opinion and a message in natural language that expresses the agent's opinion.\footnote{We use $N=32$ agents with unique personas in our experiments unless stated otherwise.} 

\textit{Communication Network.} A group of $N$ agents is connected by a communication network that determines which agents are allowed to send and receive messages from one another and the nature of the interaction between any two connected agents. We encode this network as $J \in \{-1,0,+1\}^{N\times N}$, where
a social tie $J_{ij}=+1$ means that agent $i$ is \textit{friendly (concordant)} with agent $j$, $J_{ij}=-1$ means that agent $i$ is \textit{unfriendly (discordant)} with agent $j$, and $J_{ij}=0$ means that the two agents do not communicate. These social ties model settings where an agent is more likely to trust opinions or information from certain agents and less likely for other neighbors (e.g. they might be instructed to do so by the agent owners). This formulation also includes a special case where all the edges have the same sign, so that there is no asymmetry in how agents regard each other. 
\footnote{We use symmetric $J$ with $0$ diagonal in all our experiments.} We generate the graphs $J$ from four families, which are demonstrated in Figure \ref{fig:2_Setup_Js} and described in Appendix~\ref{apdx:communication_networks}.

\paragraph{Procedure.}
All agents advance in synchronous rounds $t = 0,1,\dots,T$:

\textit{Sampling Opinions and Messages.} At each round $t$, agent $i$ expresses an answer to the question as a binary vote  $o_i(t) \sim \pi_i(\;\cdot\; | x_i(t))$ where $o_i(t) \in \{+1, -1\}$, $\pi$ is the language model, $\pi_i$ denotes conditioning on agent $i$'s persona, and $x_i(t)$ represents the messages and the question that are in the context of agent $i$ at timestep $t$. To reduce the sampling noise, we sample each agent's vote $K=5$ times and obtain $o_{i,k}(t)$ for $k = 1,\dots,K$. Then, we set $\bar{o}_i(t) = \frac{1}{K}\sum_{k=1}^K o_{i,k}(t)$ to be the opinion, so it takes one of six evenly spaced values in $[-1, 1]$. During message composition, the agent chooses a state following $s_i(t) = \text{sign}(\bar{o}_i(t))$ and samples a short message in natural language, $\omega_i(t) \sim \pi_i(\;\cdot\; | x_i(t), s_i(t) )$. The message contains its current position and provides one supporting reason.

\textit{Inbox Routing.} Each message is filed in the receiver's inbox according to the sign of their connection. If $J_{ij}=+1$ ($J_{ij}=-1$) and agent $i$ sends a message to agent $j$, the message arrives on a $+1$ ($-1$) edge and lands in the receiver's \textit{messages from friendly connections} (\textit{messages from unfriendly connections}) inbox (see Figure \ref{fig:opinion_prompt_examples} and Appendix \ref{apdx:prompts})

\textit{Opinion update.} Each agent considers its persona, the question, and its two inboxes, and re-generates its opinion to compose its message for the next turn. The updates are Markovian, \textit{i.e.},  the inboxes are refreshed at the start of each round and contain only the latest messages. Overall, influence propagates each round through the messages and opinions exchanged along the graph defined by
$J$. At $t=0$, each agent forms an opinion from its persona ($\{p_i\}_{i=1}^n$) and the question ($q$) alone.

Given the personas, the question, and the communication network, the system evolves into a terminal state ($s_1(T)$,  $s_2(T)$, \dots, $s_N(T)$), where $t=T$ is the final timestep.\footnote{We set $T=8$ in our experiments, unless stated otherwise.}
$$
(\{p_i\}_{i=1}^N,\ q,\ J)\ \longmapsto\ (s_1(t),  s_2(t), \dots, s_N(t)) \quad \text{for} \quad t = 0,1,\dots,T.
$$
We examine the behavior of this system when answering both objective and subjective questions. These two settings differ both in the questions posed and in the way personas are constructed as we described above. We provide further details on the questions, the personas, and the dataset construction for both settings in Appendix~\ref{apdx:tasks}. The backbone of the agents is a language model. We list the models we use and the embeddings we compute in Appendix~\ref{apdx:models}. For simplicity, all the communications and updates are synchronized in each round in the simulation\footnote{Notably, synchronized updates can lead to qualitative disagreements between discretized and continuous updates, \textit{e.g.}, emergence of oscillations in an energy-based system, as shown in Figure \ref{fig:1_intro}.}. We also extend the setup to asynchronous updates with similar results, which we investigate in Appendix \ref{apdx:continuous-extension}. Refer to Table \ref{tab:glossary} for a full glossary.

While our setup is simple, it captures the structure of multi-agent systems used in practice. For example, Mixture-of-Agents \citep{wang2024mixtureofagentsenhanceslargelanguage} has each agent read the outputs of all agents from the previous round, which is analogous to our procedure with a fully connected, all-friendly $J$ and depth $T$. Similarly, debate and ensemble methods correspond to other choices of $J$ and $T$ \citep{Swanson2024.11.11.623004,  du2023improvingfactualityreasoninglanguage}, and self-consistency-style methods correspond to $J$ being the identity, with every agent seeing only its own response from the previous turn \citep{wang2023selfconsistencyimproveschainthought}. Many agentic systems also use critique mechanisms, where an agent is asked to point out weaknesses of another agent's response, sometimes with agents told to disagree on purpose \citep{liang-etal-2024-encouraging}, similar to $J_{ij} = -1$. Meanwhile, systems used for scientific discovery, such as AlphaEvolve \citep{novikov2025alphaevolvecodingagentscientific}, ShinkaEvolve \citep{lange2025shinkaevolveopenendedsampleefficientprogram}, and SimpleTES \citep{ye2026structuredscalingaidiscovery}, fit a relaxation of the same template with candidate solutions as messages and the sampling policy as a dynamic $J$. 
